\documentclass{ctexart}

% 支持从项目根目录或当前文件目录编译。
\makeatletter
\def\input@path{{../}{../../}}
\makeatother
\usepackage{researchnotes}

\title{递推根式数列的极限}
\author{}
\date{}

\begin{document}
\maketitle

\section{题目与思路}

\begin{example}[原题 1.5.1]
已知实数列 $\{a_n\}_{n\ge1}$ 满足
\begin{equation}\label{eq:radical-recurrence}
  a_1=\sqrt{6},\qquad
  a_n=\sqrt{6+a_{n-1}}\quad(n=2,3,\ldots).
\end{equation}
试证 $\displaystyle\lim_{n\to\infty}a_n$ 存在，并求其值。
\end{example}

\begin{analyze}
若数列收敛于 $L$，由递推关系可预期 $L=\sqrt{6+L}$，
解得非负的候选值 $L=3$。但这一计算以收敛为前提，不能代替收敛性的证明。
因此，先证明各项小于 $3$ 且数列单调递增，再运用
\keyterm{单调有界定理}（monotone convergence theorem）：
单调递增且有上界的实数列必收敛。最后对递推关系取极限，确定极限值。
\end{analyze}

\section{解题过程}

\begin{proof}
\textbf{第一步：证明数列各项均有定义，且 $0<a_n<3$。}

当 $n=1$ 时，$0<a_1=\sqrt{6}<3$。
假设对某个整数 $n\ge1$，$a_n$ 已有定义且 $0<a_n<3$，则
\[
  6<6+a_n<9.
\]
因此 $a_{n+1}=\sqrt{6+a_n}$ 有定义，且
$0<a_{n+1}<3$。由数学归纳法，对所有整数 $n\ge1$，均有
\begin{equation}\label{eq:radical-bounds}
  0<a_n<3.
\end{equation}
特别地，$3$ 是数列的一个上界。

\textbf{第二步：证明数列严格单调递增。}

对任意整数 $n\ge1$，将相邻两项之差有理化，得
\begin{align*}
  a_{n+1}-a_n
  &=\sqrt{6+a_n}-a_n\\
  &=\frac{6+a_n-a_n^2}{\sqrt{6+a_n}+a_n}\\
  &=\frac{(3-a_n)(a_n+2)}{\sqrt{6+a_n}+a_n}>0.
\end{align*}
这里，由 \eqref{eq:radical-bounds}，分子中的两个因子以及分母均为正数。
所以 $a_{n+1}>a_n$，数列严格单调递增。

\textbf{第三步：证明极限存在，并求极限值。}

由前两步，数列单调递增且有上界。根据单调有界定理，存在有限实数 $L$，使得
\[
  \lim_{n\to\infty}a_n=L,\qquad \sqrt{6}\le L\le3.
\]
移去首项不改变数列的极限，故 $a_{n+1}\to L$。
又因平方根函数在 $6+L>0$ 处连续，对
$a_{n+1}=\sqrt{6+a_n}$ 两边取极限，得到
\[
  L=\sqrt{6+L}.
\]
两边平方并整理，有
\[
  L^2-L-6=0,\qquad (L-3)(L+2)=0.
\]
由于 $L\ge\sqrt{6}>0$，排除 $L=-2$，于是
\begin{equation}\label{eq:radical-limit}
  \boxed{\displaystyle\lim_{n\to\infty}a_n=3}.
\end{equation}
\end{proof}

\end{document}
